\scanpage{467}
\mspage{536}

\unsure{invariant} under $T_0$. \deleted{\unsure{We}}

The operation of $\mathcal M_{\rho,\sigma}^{[0]}$ on $SG$ takes place
through $\Theta$, only recently regarded as an epimorphism
\diagram{\begin{tikzcd}[ampersand replacement=\&]
\mathcal M_{\rho,\sigma}^{[0]}
\arrow[r,"\Theta|\mathcal M_{\rho,\sigma}^{[0]}"]\&B'_0
\end{tikzcd}};
this operation being specified, we can introduce
\margin{for a more \unsure{transparent} definition, cf. pages 538', 539}
\[
(17)\qquad\deleted{\unsure{\mathcal S}}
S\mathcal M_{\rho,\sigma}\overset{\text{definition}}{=}
\begin{gathered}
\mathcal M_{\rho,\sigma}^{[0]}.SG\\
\text{product }\frac12\text{ direct}
\end{gathered}
\]
in which $SG$ is invariant, and $\mathcal M_{\rho,\sigma}^{[0]}$ is identified
with a subgroup of $S\mathcal M_{\rho,\sigma}$. Since
$\mathcal M_{\rho,\sigma}^{(0)}\subset\mathcal M_{\rho,\sigma}^{[0]}$,
we have a canonical inclusion
\[
(18)\qquad
\underbrace{\mathcal M_{\rho,\sigma}^{(0)}:SG}_{\unsure{\simeq\mathcal M_{\rho,\sigma}}}
\subset
\underbrace{\mathcal M_{\rho,\sigma}^{[0]}.SG}_{=S\mathcal M_{\rho,\sigma}}
\]
whose image is \deleted{\unsure{denoted} $S\mathcal M_{\rho,\sigma}^{(0)}$}
identified with $\mathcal M_{\rho,\sigma}$.
\deleted{That said, the subgroup $L_0$ of $\mathcal M_{\rho,\sigma}^{[0]}$
cf. (16) \unsure{operates} trivial\textsuperscript{ly} on $SG$,
i.e. $\Theta\inserted{|\mathcal M_{\rho,\sigma}^{[0]}}$,
which is indeed trivial on $L_0$}

\deleted{\unsure{Let us return|Let us resume}} I am beginning \unsure{again} to get lost in it, I am going to hang on
--- for an understanding of the ``geometric'' meaning of
$S\mathcal M_{\rho,\sigma}$, $\mathcal M_{\rho,\sigma}^{[0]}$, to the notations
of page 529, which deserves to be \unsure{made explicit|made explicit} by a

\mspage{466}

\scanpage{468}

\uline{Proposition}\quad Let us consider the homomorphism
\diagram{\begin{tikzcd}[ampersand replacement=\&,column sep=large,row sep=large]
(19)\quad\mathcal M_{\rho,\sigma}\arrow[r]
\&G\times T_\rho\times N'_\sigma\times\mathbf G_m\times
\unsure{\mu_2^2|\mu_2}
\deleted{\unsure{(\mathbf G_m^3\times\mathbf G_m\times\mu_2)}}\\
x\arrow[r,mapsto]
\&(\Theta_G(x),\Theta_\rho(x),\Theta_\sigma(x),\chi(x),\varepsilon_2(x),\varepsilon_3(x))
\end{tikzcd}}
This is a monomorphism, whose image
consists of the \deleted{triples}\inserted{\unsure{systems}}
$(\underline U,\underline a,\underline b,\inserted{\mu,\varepsilon,\varepsilon'})$
satisfying
\[
\deleted{\unsure{\mu=|\text{for}}\,
\unsure{\det_G|\delta_G}(\underline U)
=\delta_\rho(\underline b)=\varepsilon\det}
\]
\[
(20)\qquad\left\{\begin{aligned}
\mu&=\delta_G(\underline U)=\delta_\sigma(\underline b)
=\varepsilon'\delta_\rho(\underline a)\\
\varepsilon&=\varepsilon_\sigma(\underline b)\\[1ex]
\end{aligned}\right.
\]

\uline{Corollary}. The morphism.
\diagram{\begin{tikzcd}[ampersand replacement=\&,column sep=large,row sep=large]
{}\arrow[d,no head,"(21)"']\&\mathcal M_{\rho,\sigma}\arrow[r]
\&G\times T_\rho\times N'_\sigma\\
{}\&x\arrow[r,mapsto]\&\Theta_G(x),\Theta_\rho(x),\Theta_\sigma(x)
\end{tikzcd}}
is a monomorphism, the image consists of the triples
$(\underline U,\underline a,\underline b)$ satisfying
\margin{\uline{N.B.} we have
$\left\{\begin{aligned}
\varepsilon&=\varepsilon_\sigma(\underline b)\\
\varepsilon'&=\delta_\rho(\underline a)/\delta_\sigma(\underline b)\\
\mu&=\delta_\sigma(\underline b)=\delta_G(\underline U)
\end{aligned}\right.$}
\[
(22)\qquad\left\{\begin{aligned}
\delta_G(\underline U)&=\delta_\sigma(\underline b),\\
[\delta_\rho(\underline a)/\delta_\sigma(\underline b)]^2&=1.
\end{aligned}\right.
\]

\uline{Proof of the proposition.}\quad We have
\deleted{$\unsure{\operatorname{Ker}(\Theta_G)|\mu_G}$}
\[
\operatorname{Ker}(\Theta_G,\deleted{\unsure{\chi|\xi_2}},\varepsilon_2,\varepsilon_3)
=SR_{\rho,\sigma}^{\unsure{\circ|0}}\simeq ST_\rho\times ST_\sigma
\]
and $\Theta_\rho,\Theta_\sigma$ \deleted{\unsure{thus|of the}}
on $SR_{\rho,\sigma}^{\unsure{\circ|0}}$ are identified with the two projections,
hence \deleted{$\unsure{\Theta}$}
\[
\operatorname{Ker}\bigl((\Theta_\rho,\Theta_\sigma)|SR_{\rho,\sigma}^{\unsure{\circ|0}}\bigr)
=\{1\},\quad\text{hence}
\]
\[
\operatorname{Ker}(\Theta_G,\chi,\varepsilon_2,\varepsilon_3,\Theta_\rho,\Theta_\sigma)
=(1),
\]

\mspage{467}

\scanpage{469}
\mspage{537}

hence a monomorphism. For the epimorphism, we are reduced to the equivalent form of the Corollary. --- We may assume that we have a $\lambda$ such that
\[
\lambda^2=\mu\overset{\mathrm{dfn}}{=}\det\underline U
\]
we want to write
\[
\underline U=\Theta_{\widetilde{\mathcal M}}
\left(\underset{\underset{R_{\rho\sigma}}{\rotatebox{90}{$\in$}}}{r},
\underset{\underset{SG}{\rotatebox{90}{$\in$}}}{U_0}\right)
=\kappa(r)U_0
\]
\margin{N.B. $\Theta_{\widetilde{\mathcal M}}$ the canonical homomorphism
\diagram{\begin{tikzcd}[ampersand replacement=\&]
\widetilde{\mathcal M}_{\rho,\sigma}\overset{\mathrm{dfn}}{=}R_{\rho,\sigma}\times SG \arrow[r] \& G
\end{tikzcd}}
which I \unsure{had|wrote|write} $\Theta_G$.}
with $U_0\in\Gamma(SG)$, i.e. $\kappa(r)^2=\det\underline U$ i.e. $\kappa(r)^2=\mu$,
we shall take $r$ with \deleted{$\unsure{\kappa}$}
\[
\kappa(r)=\lambda,
\]
\deleted{\unsure{then}} we therefore put
\[
\underline U_0=\lambda^{-1}U,\quad
\underline a_0=\lambda^{-1}\underline a,\quad
\underline b_0=\lambda^{-1}b
\]
\deleted{\unsure{and we are}} \inserted{and we are} reduced to seeing that $\underline U_0,\underline a_0,\underline b_0$ are in the image of $\mathcal M_{\rho,\sigma}$ (in \unsure{being|removing} ``\unsure{corrected ones}'' $(\underline U,\underline a,\underline b)$ by the image \deleted{\unsure{in the}} $G\times T_\rho\times\unsure{T_\sigma}$ of the section $\lambda$ of $G_m\hookrightarrow R_{\rho,\sigma}$). We are reduced to seeing that $(\underline a_0,\underline b_0)$ is in the image of
\diagram{\begin{tikzcd}[ampersand replacement=\&]
R_{\rho,\sigma}\arrow[r] \& \deleted{\unsure{T}}T_\rho\times N'_\sigma.
\end{tikzcd}}
But now $\underline a_0$ is in $S'T_\rho$, $\underline b_0$ in $SN'_\sigma$, and precisely $R_{\rho,\sigma}\simeq S'T_\rho\times SN'_\sigma$, \unsure{we have won}\ldots

\uline{Remark}. We cannot replace $G$ by $PG$,
i.e.
\diagram{\begin{tikzcd}[ampersand replacement=\&]
\mathcal M_{\rho,\sigma}\arrow[r] \& PG\times T_\rho\times N'_\sigma
\end{tikzcd}}
(where we can still add $G_m\times\mu_2\times\mu_2$ if we want!) is \uline{not} a monomorphism --- its kernel is precisely the subgroup $\mu_2$ of $\widetilde{\unsure{G|\mathcal C}}_{\rho,\sigma}^{+}=SG$.

\mspage{468}

\scanpage{470}

Let us consider, now that the quantities $\underline U,\underline a,\underline b$ $(\inserted{and}\ \mu,\varepsilon,\varepsilon')$ of page~529 (152) are understood, the quantities (153) of loc. cit. $\underline u,f,\alpha,\beta,\alpha_0,\beta_0$ ($g,h\ \xi,\eta$ do not appear there, besides\ldots) They are determined in terms of the $\underline U,\underline a,\underline b$, as soon as we know either $\underline u$ or $f$, by the formulas (154), $\underline u$ and $f$ being related by the single relation
\[
\underline U=f^{-1}\underline u \tag{23}
\]
with in addition on $f$ the condition
\[
f\in\Gamma\underbrace{SG}_{\widetilde{\unsure{G|\mathcal C}}_{\rho,\sigma}^{+}}^{\unsure{\sim}},\tag{24}
\]
plus a condition on $u$ which, in loc. cit., was written
\[
\underline u\in\Gamma T_0,\tag{25}
\]
and was suggested by the calculations of \S46 $\overline{\underline{\mathrm V}}$ (taken up again later in \S48 $\overline{\underline{\mathrm{II}}}$, $\overline{\underline{\mathrm{III}}}$ etc). This condition ensures the \uline{uniqueness} of $f,\underline u$ \unsure{satisfying|satisfying} (23) --- that is, from the $\mathrm p^{\mathrm t}$ of view of calculation, its main virtue! --- but in certain contexts it is \unsure{preferable} to replace it by the more \unsure{broad} condition
\[
\underline u\in\Gamma B'_0
\qquad\text{where}\quad
\begin{gathered}
B'_0=T_0.L_0.\\
\text{(product }\tfrac12\text{ direct)}
\end{gathered}\tag{26}
\]

\mspage{469}

\scanpage{471}
\mspage{538}

We no longer have \deleted{any \unsure{result|results}} \uline{uniqueness} for writing $\underline U$ in the form (23) --- on the other hand we have an \uline{existence} theorem in a simple form:

If we \uuline{start} with $(\alpha,\beta,f)$\deleted{$\unsure{\in}$}, sections of $G$ such that
\[
(\det\alpha)^2=\det\beta=\det f=1,\tag{27}
\]
for \deleted{us to be able to write} \inserted{this triple} and as a \unsure{result|consequence} the quantities deduced from it: \deleted{$\alpha_0,\beta_0,$}
\deleted{$\xi,\eta,g,h$}
\[
\left\{
\begin{aligned}
\alpha_0&=f^{-1}\alpha,&\quad \beta_0&=f^{-1}\unsure{\varepsilon}\beta\\
\xi&=\alpha\rho(\alpha)^{-1},&\quad \eta'&=\beta\sigma(\beta)^{-1}\\
g&=\alpha_0\rho(\alpha_0)^{-1},&\quad h&=\beta_0\sigma(\beta_0)^{-1}
\end{aligned}
\right.\tag{28}
\]
to come from a system $(x,\deleted{\unsure{f}}\inserted{\underline u})$,
with $x\in\Gamma\mathcal M_{\rho,\sigma}$,
$\deleted{\unsure{f}}\inserted{\underline u}\in\Gamma\deleted{\unsure{SG|T_0}}\inserted{B'_0}$ such that
\[
f\overset{\mathrm{dfn}}{=}\underline u\,\underline U^{-1}
\deleted{\unsure{\underline u|f}}
\in\Gamma\deleted{\unsure{B'_0}}\inserted{SG},
\quad\text{i.e.}\quad\det\underline u=\det\underline U,
\tag{29}
\]
by the formulas
\[
\alpha=\underline u\,\underline a^{-1},\quad
\beta=\underline u\,\underline b^{-1},\quad
f=\underline u\,\underline U^{-1},\tag{30}
\]
it is necessary and sufficient that we have the relation (which does not involve $f$, only $\alpha,\beta$):
\[
\alpha\rho(\alpha)^{-1}=\varepsilon\beta\sigma(\beta)^{-1}\varepsilon_1(\mu-1)
\qquad
\begin{gathered}
\text{for : }\varepsilon\in\Gamma\mu_2\\
\mu\in\Gamma G_m\\
\text{suitable.}
\end{gathered}\tag{31}
\]

\uline{(N.B.} $\unsure{\varepsilon|\varepsilon_1},\mu$ are uniquely determined (\unsure{thus} in terms of $\alpha\ \beta$).
\inserted{Then the pair $(x,\underline u)$ is unique. For $f=1$, (i.e. $\underline u=\underline U$),}
We can say \deleted{that} also that we have a $1$-$1$ correspondence between the \inserted{the set \deleted{\illegible{1}} sections \unsure{of the|of the}} \unsure{subgroups|subgroup} $\mathcal M_{\rho,\sigma}[0]$ (not $\mathcal M_{\rho,\sigma}(0)$!) of $\mathcal M_{\rho,\sigma}$, and the set

\mspage{470}

\scanpage{472}

of pairs $\alpha,\beta$ satisfying the two relations
\[
(\det\alpha)^2=\det\beta=1\qquad\text{of (27)}
\]
\uline{and} the loop relation. For this, we have \deleted{\unsure{besides}} no use for the group $S\mathcal M_{\rho,\sigma}$.

But returning \unsure{to the case} of arbitrary $f$, to interpret the \unsure{algebraic} \unsure{sheaf} on the triples $(\alpha,\beta,f)$, \inserted{satisfying (27)(31),} \inserted{with the structure of} a strange group on it \inserted{\unsure{via} cocycle relation,} (to which I shall have to return in \inserted{more} \unsure{geometric} terms\ldots) and because the three quantities $\alpha,\beta,f$ mix intimately --- there is reason to introduce $S\mathcal M_{\rho,\sigma}$ --- which will be isomorphic \unsure{to the|to} scheme (in groups) of the common solutions of (27)(31).

From this point of view, $S\mathcal M_{\rho,\sigma}$ is derived from $\unsure{\rho}$ as \deleted{\unsure{described by triples}} \inserted{scheme}
\deleted{$(x,\underline u,f)$ \unsure{of}.}
\[
\begin{aligned}
S\mathcal M_{\rho,\sigma}&=\text{scheme of triples }(x,\underline u,f)\\
&\text{with }\left\{
\begin{aligned}
x\ \deleted{\unsure{\in}}&\text{ in }\mathcal M_{\rho,\sigma},\\
\deleted{\unsure{f}}\inserted{\underline u}&\text{ in }B'_0\\
f&\text{ in }SG
\end{aligned}
\right.
\qquad\left\{
\begin{aligned}
\underline U&=f^{-1}\underline u\\
\Vert\,\text{definition}&\\
\Theta_G(x)&
\end{aligned}\right.
\end{aligned}\tag{32}
\]

Thus
\diagram{\begin{tikzcd}[ampersand replacement=\&,column sep=large,row sep=small]
{(33)\quad S\mathcal M_{\rho,\sigma}}
\arrow[r,hook]
\arrow[d,phantom,"{\unsure{\rotatebox{90}{$\simeq$}}}"]
\& {\mathcal M_{\rho,\sigma}\times B'_0\quad\unsure{\leftarrow}}\\
{\left\{x,\underline u\ \middle|\ \unsure{\chi|\gamma}(x)=
\underbrace{\unsure{\delta_G}(\underline u)}_{\begin{gathered}\Vert\,\mathrm{dfn}\\\det\underline u\end{gathered}}\right\}} \&
\end{tikzcd}}
or again, more simply

\mspage{471}